% Options for packages loaded elsewhere
% Options for packages loaded elsewhere
\PassOptionsToPackage{unicode}{hyperref}
\PassOptionsToPackage{hyphens}{url}
\PassOptionsToPackage{dvipsnames,svgnames,x11names}{xcolor}
%
\documentclass[
  number,
  preprint]{elsarticle}
\usepackage{xcolor}
\usepackage{amsmath,amssymb}
\setcounter{secnumdepth}{5}
\usepackage{iftex}
\ifPDFTeX
  \usepackage[T1]{fontenc}
  \usepackage[utf8]{inputenc}
  \usepackage{textcomp} % provide euro and other symbols
\else % if luatex or xetex
  \usepackage{unicode-math} % this also loads fontspec
  \defaultfontfeatures{Scale=MatchLowercase}
  \defaultfontfeatures[\rmfamily]{Ligatures=TeX,Scale=1}
\fi
\usepackage{lmodern}
\ifPDFTeX\else
  % xetex/luatex font selection
\fi
% Use upquote if available, for straight quotes in verbatim environments
\IfFileExists{upquote.sty}{\usepackage{upquote}}{}
\IfFileExists{microtype.sty}{% use microtype if available
  \usepackage[]{microtype}
  \UseMicrotypeSet[protrusion]{basicmath} % disable protrusion for tt fonts
}{}
\makeatletter
\@ifundefined{KOMAClassName}{% if non-KOMA class
  \IfFileExists{parskip.sty}{%
    \usepackage{parskip}
  }{% else
    \setlength{\parindent}{0pt}
    \setlength{\parskip}{6pt plus 2pt minus 1pt}}
}{% if KOMA class
  \KOMAoptions{parskip=half}}
\makeatother
% Make \paragraph and \subparagraph free-standing
\makeatletter
\ifx\paragraph\undefined\else
  \let\oldparagraph\paragraph
  \renewcommand{\paragraph}{
    \@ifstar
      \xxxParagraphStar
      \xxxParagraphNoStar
  }
  \newcommand{\xxxParagraphStar}[1]{\oldparagraph*{#1}\mbox{}}
  \newcommand{\xxxParagraphNoStar}[1]{\oldparagraph{#1}\mbox{}}
\fi
\ifx\subparagraph\undefined\else
  \let\oldsubparagraph\subparagraph
  \renewcommand{\subparagraph}{
    \@ifstar
      \xxxSubParagraphStar
      \xxxSubParagraphNoStar
  }
  \newcommand{\xxxSubParagraphStar}[1]{\oldsubparagraph*{#1}\mbox{}}
  \newcommand{\xxxSubParagraphNoStar}[1]{\oldsubparagraph{#1}\mbox{}}
\fi
\makeatother


\usepackage{longtable,booktabs,array}
\newcounter{none} % for unnumbered tables
\usepackage{calc} % for calculating minipage widths
% Correct order of tables after \paragraph or \subparagraph
\usepackage{etoolbox}
\makeatletter
\patchcmd\longtable{\par}{\if@noskipsec\mbox{}\fi\par}{}{}
\makeatother
% Allow footnotes in longtable head/foot
\IfFileExists{footnotehyper.sty}{\usepackage{footnotehyper}}{\usepackage{footnote}}
\makesavenoteenv{longtable}
\usepackage{graphicx}
\makeatletter
\newsavebox\pandoc@box
\newcommand*\pandocbounded[1]{% scales image to fit in text height/width
  \sbox\pandoc@box{#1}%
  \Gscale@div\@tempa{\textheight}{\dimexpr\ht\pandoc@box+\dp\pandoc@box\relax}%
  \Gscale@div\@tempb{\linewidth}{\wd\pandoc@box}%
  \ifdim\@tempb\p@<\@tempa\p@\let\@tempa\@tempb\fi% select the smaller of both
  \ifdim\@tempa\p@<\p@\scalebox{\@tempa}{\usebox\pandoc@box}%
  \else\usebox{\pandoc@box}%
  \fi%
}
% Set default figure placement to htbp
\def\fps@figure{htbp}
\makeatother





\setlength{\emergencystretch}{3em} % prevent overfull lines

\providecommand{\tightlist}{%
  \setlength{\itemsep}{0pt}\setlength{\parskip}{0pt}}



 
\usepackage[]{natbib}
\bibliographystyle{elsarticle-num-names}


\makeatletter
\@ifpackageloaded{caption}{}{\usepackage{caption}}
\AtBeginDocument{%
\ifdefined\contentsname
  \renewcommand*\contentsname{Table of contents}
\else
  \newcommand\contentsname{Table of contents}
\fi
\ifdefined\listfigurename
  \renewcommand*\listfigurename{List of Figures}
\else
  \newcommand\listfigurename{List of Figures}
\fi
\ifdefined\listtablename
  \renewcommand*\listtablename{List of Tables}
\else
  \newcommand\listtablename{List of Tables}
\fi
\ifdefined\figurename
  \renewcommand*\figurename{Figure}
\else
  \newcommand\figurename{Figure}
\fi
\ifdefined\tablename
  \renewcommand*\tablename{Table}
\else
  \newcommand\tablename{Table}
\fi
}
\@ifpackageloaded{float}{}{\usepackage{float}}
\floatstyle{ruled}
\@ifundefined{c@chapter}{\newfloat{codelisting}{h}{lop}}{\newfloat{codelisting}{h}{lop}[chapter]}
\floatname{codelisting}{Listing}
\newcommand*\listoflistings{\listof{codelisting}{List of Listings}}
\makeatother
\makeatletter
\usepackage{pdflscape}
\makeatother
\makeatletter
\makeatother
\makeatletter
\@ifpackageloaded{caption}{}{\usepackage{caption}}
\@ifpackageloaded{subcaption}{}{\usepackage{subcaption}}
\makeatother
\journal{SSRN}
\usepackage{bookmark}
\IfFileExists{xurl.sty}{\usepackage{xurl}}{} % add URL line breaks if available
\urlstyle{same}
% fallback for those not using the hyperref driver hyperxmp:
\makeatletter
\@ifundefined{xmpquote}{\newcommand{\xmpquote}[1]{#1}}{}
\makeatother
\hypersetup{
  pdftitle={Whose experience counts? A narrative case review of citizen science for public health research},
  pdfkeywords={\xmpquote{Citizen Science}},
  colorlinks=true,
  linkcolor={blue},
  filecolor={Maroon},
  citecolor={Blue},
  urlcolor={Blue},
  pdfcreator={LaTeX via pandoc}}


\setlength{\parindent}{6pt}
\begin{document}

\begin{frontmatter}
\title{Whose experience counts? A narrative case review of citizen
science for public health research}
\author[1,2]{Aly Lamuri%
\corref{cor1}%
}
 \ead{aly.lamuri@ui.ac.id} 
\author[3]{Feri Setiawan%
%
}

\author[1,5]{Haykal Hafidzul Arifin%
%
}

\author[4]{Firas Farisi Alkaff%
%
}

\author[2]{Novika Purnama Sari%
%
}

\author[3]{Nur Afny Catur Andryani%
%
}


\affiliation[1]{organization={University of
Indonesia},city={Depok},country={ID},countrysep={,},postcodesep={}}
\affiliation[2]{organization={University of
Groningen},city={Groningen},country={Netherlands},countrysep={,},postcodesep={}}
\affiliation[3]{organization={Binus
University},city={Jakarta},country={ID},countrysep={,},postcodesep={}}
\affiliation[4]{organization={University Medical Center
Groningen},city={Groningen},country={NL},countrysep={,},postcodesep={}}
\affiliation[5]{organization={PT Atensi
Kini},city={Jakarta},country={ID},countrysep={,},postcodesep={}}

\cortext[cor1]{Corresponding author}






        
\begin{abstract}
\textbf{Background.} Stress, depression, anxiety and burnout are common,
but most people who experience them never reach services, so their
experience rarely becomes data. Citizen science, in which members of the
public take part in research as partners rather than only as subjects,
has been proposed as a way to close this gap. Evidence on how it works
in mental health comes almost entirely from high-income countries and
rarely includes projects that failed.\\
\textbf{Aim.} To describe how citizens took part across the stages of
research in a set of successful and failed cases, what changed because
they took part, and what these cases teach for building citizen-centred
mental health research in Indonesia.\\
\textbf{Methods.} A narrative review of eight purposively selected
cases: four mental-health citizen science projects, two Indonesian
citizen health-data initiatives, and two failures involving
mental-health data. Each case was scored on an 18-item framework
covering what changed (three items), where citizens took part (six
items), how the project was run (seven items), and durability and harm
(two items), using a four-level scale (met, partly met, not met, not
reported). Sources were peer-reviewed articles, preprints and
organisational reports, with news used only where nothing else existed.
\end{abstract}





\begin{keyword}
    
    Citizen Science
\end{keyword}
\end{frontmatter}
    

\section{Introduction}\label{sec-introduction}

Stress, low mood, worry and exhaustion are among the most widely shared
human experiences. People talk about them with friends, colleagues and
online communities long before, and often without ever, receiving a
diagnosis. When many people describe the same experience in similar
words, that shared description becomes a signal worth studying. Long
COVID is the clearest recent example: patients named and documented the
condition through online communities, and within months it moved from
social media into clinical and policy channels, before most formal
studies had begun \citep{callard2021}. This kind of knowledge is easily
dismissed. Leaving people out of how knowledge about their own
experience is produced has been described as a form of epistemic
injustice, meaning a harm done to someone in their capacity as a knower
\citep{okoroji2023}. The starting point of this review is that the
collective experience of distress is evidence waiting to be given
structure.

The scale of that experience is large, and most of it never reaches
health services. Globally, the COVID-19 pandemic was associated with an
estimated 53.2 million additional cases of major depressive disorder (an
increase of 27.6\%) and 76.2 million additional cases of anxiety
disorders (an increase of 25.6\%) in 2020 alone \citep{covidmental2021}.
In Indonesia, analysis of the 2018 national basic health survey found
depression in 5.1\% of adolescents aged 15 to 17 and 5.6\% of young
adults aged 18 to 24 \citep{suryaputri2022}. The first national
adolescent mental health survey found that about one in twenty
adolescents aged 10 to 17 had a mental disorder in the previous year,
while only about one in fifty had used any service for emotional or
behavioural problems \citep{inamhs2022}. Burnout is similarly
widespread. In an online assessment of 4,338 full-time employees in
Indonesia, Malaysia, the Philippines and Singapore in 2022, 62.9\%
reported high or very high burnout \citep{abdulaziz2024}, and a
meta-analysis estimated burnout in 39\% of the public health workforce
worldwide \citep{nagarajan2024}. In Indonesia, a national survey of
3,629 healthcare workers likewise found burnout syndrome in 37.5\%
\citep{lamuri2023}. Measures of resilience, the capacity to cope and
recover, are rarely collected alongside these measures of distress, so
what helps people stay well is even less visible than what harms them.
Because most distress does not reach services, data drawn from services
describe only a small, unrepresentative part of it.

Conventional research is poorly placed to fill this gap. Questions are
usually chosen by researchers or institutions, and members of the public
take part mainly by filling in forms. In a regional study that included
Indonesia, community members said they did not know how the information
they gave was used \citep{kumar2025}. This matters for the quality of
research as well as for fairness. In a longitudinal study of mental
health research in England, studies with greater patient involvement
were more likely to reach their recruitment targets, although the
association does not prove cause \citep{ennis2013}. Involvement can also
go wrong: when lived-experience roles are filled by a small,
well-connected group, or when people are asked to share painful
experience without fair support and pay, involvement risks ``elite
capture'' and ``epistemic exploitation'' \citep{okoroji2023}. Research
done on people, rather than with them, therefore risks asking the wrong
questions, recruiting the wrong people and losing their trust.

Citizen science offers a different arrangement. In this paper, citizen
science means research in which members of the public take part as
partners: choosing questions, designing methods, collecting or
interpreting data, and sharing results. Participation ranges from
contributing data (crowdsourcing) to defining the problem and leading
the research (extreme citizen science) \citep{haklay2013}, and the
European Citizen Science Association's ten principles describe what good
practice looks like \citep{robinson2018}. In mental health specifically,
a systematic review found that citizen scientists with lived experience
reported empowerment, social connection and improved wellbeing and
self-esteem, and that the main cost was the time and energy the work
required \citep{todowede2023}. Citizen science can also reach scale: a
Dutch crowdsourcing study recruited 12,503 members of the public to
report both symptoms and personal strengths, and gave each participant
personalised feedback in return \citep{vanderkrieke2016}. Citizen
science therefore has the potential to make research on distress and
resilience more relevant, larger in scale and more reciprocal.

Indonesia appears ready for this approach but is not yet familiar with
it. In a nine-country study, 82.2\% of Indonesian respondents had never
heard of citizen science, yet 86.3\% owned a smartphone and 91.6\% had
internet access; community support, including mental support, stood out
as a way Indonesians wanted to contribute \citep{tan2024}. During the
COVID-19 pandemic, grassroots groups began collecting and publishing
health data themselves because official information was inadequate
\citep{nadzir2020}, and one such group documented health-worker deaths
that no official source was counting \citep{ekawati2022}. The conditions
for phone-based, citizen-centred mental health research exist; the
practice does not yet.

The evidence base for that practice has three gaps. First, all nine
studies in the systematic review of citizen science in mental health
were conducted in high-income countries \citep{todowede2023}. Second,
the same review found that studies rarely reported feedback to citizen
scientists and none reported making their data publicly available, so it
is hard to judge how projects were actually run \citep{todowede2023}.
Third, failures are rarely written up, which biases lessons toward
success. This review addresses these gaps by comparing successful and
failed cases on a single scoring framework. It asks three questions: (1)
What does citizen involvement look like at each stage of research, from
choosing the question to acting on findings? (2) Which initiatives in
Indonesia and elsewhere have done this, and what happened? (3) What do
these cases teach for building citizen-centred mental health research?
It is the first paper in a series: its lessons inform a set of
principles for research-ready health data (Paper 2), an assessment of an
open-source research platform against those principles (Paper 3), and a
baseline study of depression, anxiety, burnout, resilience and
functioning among adults in three Indonesian cities (Paper 4).

\section{Methods}\label{sec-methods}

\subsection{Design}\label{design}

We conducted a narrative review of multiple cases using publicly
available documents. We chose a narrative rather than a systematic
design because the number of relevant cases is small, much of the
evidence is grey literature (reports and documents not published in
journals), and the aim was to draw lessons rather than to estimate an
effect. The review was checked against SANRA, a six-item scale for
assessing the quality of narrative reviews \citep{baethge2019}. No
ethical approval was required because only public documents were used.

\subsection{Definitions}\label{definitions}

A \emph{citizen} was defined as anyone acting outside a paid research or
government role. Because this definition includes civil-society
professionals who volunteer their expertise, we recorded separately for
each case whether leading roles were held by lay members of the public,
by patient-researchers, by civil-society experts, or by institutions.
The \emph{unit of analysis} was the citizen-data activity rather than
the organisation as a whole. For example, in a case where an
organisation shared data from its service, we scored the data-sharing
programme, not the service.

\subsection{Sources}\label{sources}

Cases were first identified through scoping searches of the web and the
reference lists of known projects. To check for missed cases, we then
ran structured searches on 4 November 2026 in PubMed, Scopus and Europe
PMC (which also indexes preprints), Google Scholar (first 200 results
per query, ranked by relevance), the Indonesian journal index Garuda,
and two citizen-science project directories (EU-Citizen.Science and
SciStarter), in English and Bahasa Indonesia. Three query sets covered
citizen science in mental health, participatory mental-health research
in Indonesia, and controversies over the use of mental-health data.
Exact queries, filters, numbers of results and screening outcomes are
reported in Supplement S4, following the PRISMA-S extension for
reporting literature searches \citep{rethlefsen2021}. For each included
case, we then retrieved project websites, organisational reports and the
reference lists of retrieved documents. Sources were preferred in the
following order: peer-reviewed article, preprint or conference
proceedings, organisational report, project website, and news. News was
used only where no other source documented a fact, and this is noted in
the text.

\subsection{Case selection}\label{case-selection}

Cases were selected purposively. To be included, a case had to (1)
concern health, with priority for stress, depression, anxiety, burnout
or resilience; (2) give citizens a role beyond being data subjects in at
least one stage of research, or, for failure cases, involve data from
people with mental-health needs; (3) have documented outcomes, or be
clearly identified as ongoing; and (4) be described in at least two
independent sources, at least one of which was a peer-reviewed article,
preprint or organisational report. We sought a mix of Indonesian and
international cases, of successes and failures, and of citizen-led and
institution-led projects. Table~\ref{tbl-cases} lists the candidate
cases and the reason for including or excluding each.

\begin{longtable}[]{@{}
  >{\raggedright\arraybackslash}p{(\linewidth - 6\tabcolsep) * \real{0.1800}}
  >{\raggedright\arraybackslash}p{(\linewidth - 6\tabcolsep) * \real{0.1400}}
  >{\raggedright\arraybackslash}p{(\linewidth - 6\tabcolsep) * \real{0.1800}}
  >{\raggedright\arraybackslash}p{(\linewidth - 6\tabcolsep) * \real{0.5000}}@{}}
\caption[]{Candidate cases and selection
decisions}\label{tbl-cases}\tabularnewline
\toprule\noalign{}
\begin{minipage}[b]{\linewidth}\raggedright
Case
\end{minipage} & \begin{minipage}[b]{\linewidth}\raggedright
Setting
\end{minipage} & \begin{minipage}[b]{\linewidth}\raggedright
Role in review
\end{minipage} & \begin{minipage}[b]{\linewidth}\raggedright
Decision and reason
\end{minipage} \\
\midrule\noalign{}
\endfirsthead
\toprule\noalign{}
\begin{minipage}[b]{\linewidth}\raggedright
Case
\end{minipage} & \begin{minipage}[b]{\linewidth}\raggedright
Setting
\end{minipage} & \begin{minipage}[b]{\linewidth}\raggedright
Role in review
\end{minipage} & \begin{minipage}[b]{\linewidth}\raggedright
Decision and reason
\end{minipage} \\
\midrule\noalign{}
\endhead
\bottomrule\noalign{}
\endlastfoot
CoAct for Mental Health & Spain, 2020--2022 & Core mental-health case &
Included \\
HoeGekIsNL & Netherlands, mid-2010s & Core mental-health case &
Included \\
Patient-Led Research Collaborative & International, 2020 & Core case
(cognitive and psychological symptoms of Long COVID) & Included \\
Youth-GEMS & European Union and Australia, 2023 & Core mental-health
case, institution-led contrast & Included (ongoing) \\
LaporCovid-19 & Indonesia, 2020 & Indonesian context case & Included \\
Community-based surveillance (Indonesian Red Cross) & Indonesia,
2018--2020 & Indonesian context case & Included \\
Crisis Text Line data sharing & United States, until 2022 & Failure case
& Included \\
Samaritans Radar & United Kingdom, 2014--2015 & Failure case &
Included \\
KawalCOVID19 & Indonesia, 2020 & Indonesian context case & Excluded:
evidence almost entirely from news reports
\citep{widianto2020, nusantara2024} \\
Hospital mercury monitoring, Denpasar & Indonesia, 2012--2013 &
Environmental health & Excluded: outside mental-health focus;
measurements taken by NGO staff, with no documented lay participation
\citep{balifokus2013} \\
\end{longtable}

\subsection{Scoring framework}\label{scoring-framework}

Each case was scored on 18 items grouped into four domains
(Table~\ref{tbl-scoring}). The items were derived from three sources:
the stages of the research lifecycle; Haklay's typology of participation
\citep{haklay2013}; and the ten principles of citizen science of the
European Citizen Science Association \citep{robinson2018}. Principles 1,
2 and 4 of the ECSA framework are covered by Domains A and C; the
remaining principles each map to one item in Domain B. A crosswalk
between items and source frameworks is given in Supplement S1.

Each item was scored on four levels: \texttt{Y} met; \texttt{P} partly
met; \texttt{N} not met; and \texttt{NR}, not reported. A score of
\texttt{N} required evidence that the practice did not happen; where
sources were silent, the item was scored \texttt{NR}. This distinction
matters because the review relies on public documents, and an unreported
practice may still have taken place.

\begin{longtable}[]{@{}
  >{\raggedright\arraybackslash}p{(\linewidth - 4\tabcolsep) * \real{0.0800}}
  >{\raggedright\arraybackslash}p{(\linewidth - 4\tabcolsep) * \real{0.3200}}
  >{\raggedright\arraybackslash}p{(\linewidth - 4\tabcolsep) * \real{0.6000}}@{}}

\caption{\label{tbl-scoring}Scoring framework}

\tabularnewline

\toprule\noalign{}
\begin{minipage}[b]{\linewidth}\raggedright
Item
\end{minipage} & \begin{minipage}[b]{\linewidth}\raggedright
Question
\end{minipage} & \begin{minipage}[b]{\linewidth}\raggedright
\texttt{Y} Met when sources show\ldots{}
\end{minipage} \\
\midrule\noalign{}
\endhead
\bottomrule\noalign{}
\endlastfoot
\multicolumn{3}{@{}>{\raggedright\arraybackslash}p{(\linewidth - 4\tabcolsep) * \real{1.0000} + 4\tabcolsep}@{}}{%
\textbf{A. Citizens' highest documented role at each stage}} \\
A1 & Choosing the question & Citizens led or co-decided \\
A2 & Designing methods and tools & Citizens led or co-decided \\
A3 & Collecting data & Citizens led or co-decided \\
A4 & Analysing and interpreting data & Citizens led or co-decided \\
A5 & Writing and sharing results & Citizens led or co-decided \\
A6 & Acting on findings & Citizens led or co-decided \\
\multicolumn{3}{@{}>{\raggedright\arraybackslash}p{(\linewidth - 4\tabcolsep) * \real{1.0000} + 4\tabcolsep}@{}}{%
\textbf{B. How the project was run}} \\
B1 & Both citizens and science benefited (ECSA 3) & A stated benefit to
participants and a scientific output \\
B2 & Participants got results back (ECSA 5) & Results returned in a form
participants could use \\
B3 & Data and results are open (ECSA 7) & Both openly available, or
under controlled access for sensitive data \\
B4 & Contributions acknowledged in a form contributors chose (ECSA 8) &
Acknowledgement, with contributors able to choose named, collective or
anonymous credit \\
B5 & Data quality checked, limitations stated (ECSA 6) & Both
described \\
B6 & Ethics and data protection in place (ECSA 10) & Ethics review and a
consent or privacy approach \\
B7 & Project evaluated itself (ECSA 9) & A formal evaluation
published \\
\multicolumn{3}{@{}>{\raggedright\arraybackslash}p{(\linewidth - 4\tabcolsep) * \real{1.0000} + 4\tabcolsep}@{}}{%
\textbf{C. What changed because citizens took part}} \\
C1 & Citizens changed the research question & A focus or measure studied
because citizens raised it, which researchers or officials were not
studying \\
C2 & Findings came before the official or usual route & Dated sources
show results published before the official or conventional equivalent \\
C3 & Findings were trusted and used & Documented use by a government,
health service or funder \\
\multicolumn{3}{@{}>{\raggedright\arraybackslash}p{(\linewidth - 4\tabcolsep) * \real{1.0000} + 4\tabcolsep}@{}}{%
\textbf{D. Durability and harm}} \\
D1 & Activity lasted, or ended as planned & Still running, or closed
with a hand-over \\
D2 & Volunteers and participants protected from harm & Risks foreseen
and managed, no harm documented \\

\end{longtable}

\emph{Note.} Full definitions in Supplement S1.

In Domain A, \texttt{Y} means citizens led or co-decided a stage,
\texttt{P} means they contributed or were consulted, and \texttt{N}
means they had no role. Each stage was scored at the highest role
documented. People who were sampled as survey respondents or who used a
service were not counted as contributing to data collection; members of
the public who enrolled themselves in an open crowdsourcing study were
counted as contributors (\texttt{P}).

Two summaries were derived from the item scores rather than scored
separately. The \emph{level of participation} followed Haklay's
typology: crowdsourcing if citizens contributed only to data collection;
distributed intelligence if they also contributed to analysis or
interpretation; participatory science if they were consulted on or
contributed to the question or design; and extreme citizen science if
they led or co-decided the question. Cases in which citizens had no
research role were labelled ``data subjects only''. The \emph{verdict}
gave most weight to Domain C, reflecting the review's focus on what
citizen involvement changes. A case was a \textbf{success} if findings
were documented as used (C3 = \texttt{Y}) and the activity lasted or
ended as planned (D1 is \texttt{P}/\texttt{Y}); a \textbf{failure} if
findings were rejected (C3 = \texttt{N}), the activity collapsed (D1 =
\texttt{N}), or harm was not addressed (D2 = \texttt{N});
\textbf{partial} if neither applied; and \textbf{insufficient evidence}
if C3 or D1 was not reported and no failure criterion was met. These
rules were fixed before scoring.

\subsection{Coding procedure}\label{coding-procedure}

Two coders scored each case independently. The framework was first
piloted on two cases, after which definitions were refined; the
refinements are listed in Supplement S1. Agreement is reported as the
percentage of cells on which coders agreed {[}to be completed: XX\%{]};
disagreements were resolved by discussion. Every score is supported by a
citation and a one-line justification in Supplement S2. To limit bias
arising from the lead author's role in a research platform (see
Declarations), the second coder had no ties to that platform and held
the final decision on disputed cells.

\subsection{Synthesis}\label{synthesis}

For each item we counted the number of cases scored \texttt{Y}. Before
scoring, we set the thresholds for the summary: an item was \emph{done
well across cases} if it was met in at least six of eight cases, and
\emph{consistently lacking} if it was met in two or fewer. Consistently
lacking items were further divided by whether the cases mostly scored
\texttt{N} (not done) or \texttt{NR} (not reported). We compared cases
by leadership (citizen-led versus institution-led), setting (Indonesian
versus international) and verdict (success versus failure), and drew a
numbered list of lessons, each linked to the cases that support it.

\section{Results}\label{sec-results}

\subsection{Overview of cases}\label{overview-of-cases}

The eight included cases span 2014 to the present and come from Europe,
Indonesia, the United States and an international patient network
(Table~\ref{tbl-case-overview}). Four are mental-health citizen science
projects, two are Indonesian citizen health-data initiatives that arose
during the COVID-19 pandemic, and two are failures involving data from
people in mental distress. Leading roles were held by the group that
founded the initiative in two cases: patients in the Patient-Led
Research Collaborative, and, in LaporCovid-19, a civil-society coalition
of public health experts, scientists, journalists and activists that
describes itself as a citizens' coalition
\citep{ekawati2022, laporcovid2021}. Leading roles were shared between a
university and lived-experience co-researchers in one case (CoAct) and
held by institutions in the remaining five. Evidence was strongest for
the Patient-Led Research Collaborative and LaporCovid-19, each
documented in peer-reviewed articles, and weakest for Youth-GEMS, which
has not yet published results, and Samaritans Radar, which is documented
in grey literature and news.

\begin{landscape}

\begin{longtable}[]{@{}
  >{\raggedright\arraybackslash}p{(\linewidth - 12\tabcolsep) * \real{0.2000}}
  >{\raggedright\arraybackslash}p{(\linewidth - 12\tabcolsep) * \real{0.1200}}
  >{\raggedright\arraybackslash}p{(\linewidth - 12\tabcolsep) * \real{0.2400}}
  >{\raggedright\arraybackslash}p{(\linewidth - 12\tabcolsep) * \real{0.1100}}
  >{\raggedright\arraybackslash}p{(\linewidth - 12\tabcolsep) * \real{0.1300}}
  >{\raggedright\arraybackslash}p{(\linewidth - 12\tabcolsep) * \real{0.1000}}
  >{\raggedright\arraybackslash}p{(\linewidth - 12\tabcolsep) * \real{0.1000}}@{}}
\caption[]{Overview of included
cases}\label{tbl-case-overview}\tabularnewline
\toprule\noalign{}
\begin{minipage}[b]{\linewidth}\raggedright
Case
\end{minipage} & \begin{minipage}[b]{\linewidth}\raggedright
Setting
\end{minipage} & \begin{minipage}[b]{\linewidth}\raggedright
Who led
\end{minipage} & \begin{minipage}[b]{\linewidth}\raggedright
Scale
\end{minipage} & \begin{minipage}[b]{\linewidth}\raggedright
Main sources
\end{minipage} & \begin{minipage}[b]{\linewidth}\raggedright
Level of participation
\end{minipage} & \begin{minipage}[b]{\linewidth}\raggedright
Verdict
\end{minipage} \\
\midrule\noalign{}
\endfirsthead
\toprule\noalign{}
\begin{minipage}[b]{\linewidth}\raggedright
Case
\end{minipage} & \begin{minipage}[b]{\linewidth}\raggedright
Setting
\end{minipage} & \begin{minipage}[b]{\linewidth}\raggedright
Who led
\end{minipage} & \begin{minipage}[b]{\linewidth}\raggedright
Scale
\end{minipage} & \begin{minipage}[b]{\linewidth}\raggedright
Main sources
\end{minipage} & \begin{minipage}[b]{\linewidth}\raggedright
Level of participation
\end{minipage} & \begin{minipage}[b]{\linewidth}\raggedright
Verdict
\end{minipage} \\
\midrule\noalign{}
\endhead
\bottomrule\noalign{}
\endlastfoot
CoAct for Mental Health & Barcelona, 2020--2022 & University with 32
lived-experience co-researchers & 394 participants analysed &
Proceedings & Participatory & Partial \\
HoeGekIsNL & Netherlands & Academic research centre & 12,503
participants & Peer-reviewed article & Crowdsourcing & Insufficient
evidence \\
Patient-Led Research Collaborative & International, 2020-- &
Patient-researchers & 3,762 respondents & Peer-reviewed articles,
preprint & Extreme & Partial \\
Youth-GEMS & EU and Australia, 2023-- & Academic consortium & Planned
\textasciitilde2,000 & Website, conference abstract & Participatory &
Insufficient evidence (ongoing) \\
LaporCovid-19 & Indonesia, 2020-- & Civil-society coalition of experts,
scientists, journalists and activists & 1,545 deaths documented &
Peer-reviewed articles & Extreme & Partial \\
Community-based surveillance & Indonesia, 2018--2020 & Indonesian Red
Cross & 234 volunteers & Organisational report & Participatory &
Success \\
Crisis Text Line data sharing & United States, until 2022 & Non-profit
organisation & National crisis service & Peer-reviewed editorial & Data
subjects only & Failure \\
Samaritans Radar & United Kingdom, 2014--2015 & Charity & National app &
Grey literature, news & Data subjects only & Failure \\
\end{longtable}

\end{landscape}

\subsection{Case narratives}\label{case-narratives}

\subsubsection{CoAct for Mental Health}\label{coact-for-mental-health}

CoAct for Mental Health was part of a European citizen social science
project funded under Horizon 2020. In 2020, an open call recruited
people with lived experience of mental health problems and family
members as co-researchers, recognised as experts in the field
\citep{casanovas2026}. The research focused on social support networks,
a need identified by the communities themselves. Co-researchers took
part in formulating the research questions, co-designed a chatbot on the
messaging app Telegram, and wrote the micro-stories of social support
that the chatbot shared with participants. The chatbot reached up to 900
people; the analysed dataset contained 11,002 interactions from 394
participants responding to 129 stories. Analysis was carried out by the
research group and interpreted collectively with the co-researchers and
a wider coalition that included representatives of the mental health
system. In November 2022, a community assembly approved 14
recommendations, which were presented to Barcelona City Council and to
the regional mental health pact, and published as policy briefs
\citep{casanovas2026}. At the assembly, a city councillor said she
wanted all validated recommendations included in the new municipal
mental health plan, and the director of the National Mental Health Pact
said the Pact would incorporate the proposed actions
\citep{ub2022assembly, mitats2022}. We found no later source tracing
whether the recommendations were adopted; the city's announcement of its
Mental Health Plan 2023--2030 does not mention the project
\citep{bcn2023plan}. The chatbot allowed anonymous participation on a
sensitive topic, and the authors note the need for emotional safety. The
project won the 2024 European Union Prize for Citizen Science.

\emph{Lesson:} shared governance with people with lived experience
produces co-owned questions and policy proposals, but whether those
proposals are used depends on commitments from decision-makers secured
before the project ends, and on a way to track whether those commitments
are kept.

\subsubsection{HoeGekIsNL}\label{hoegekisnl}

HoeGekIsNL (``How crazy is the Netherlands?''; published in English as
HowNutsAreTheDutch) is a national crowdsourcing study initiated by an
academic research centre in Groningen to map the mental strengths and
vulnerabilities of the Dutch population
\citep{vanderkrieke2016, mindwise2016}. Members of the public enrolled
themselves through a website and completed questionnaires on both
psychiatric symptoms and mental strengths and resources; a subgroup also
reported on their smartphones three times a day for 30 days. In total,
12,503 people took part \citep{vanderkrieke2016}. Participants were not
paid. Instead, they received personalised feedback and could compare
themselves with other participants; those who completed most of the
diary received an automatically generated report on patterns in their
emotions and behaviour, and on how their strengths and vulnerabilities
interacted \citep{mindwise2016}. We found no documentation of public
involvement in choosing the questions or designing the study. Although
participation was intended for adults, the study team reported that more
than 200 minors took part \citep{mindwise2016}, which illustrates a
safeguarding challenge of open online enrolment. At the time of writing,
the study website still invited new participants and also offered a
COVID-19 questionnaire \citep{hoegekis_web}.

\emph{Lesson:} personal feedback can be both the incentive that brings
people in at scale and the means of returning results to them. Measuring
strengths alongside symptoms is feasible at population scale.

\subsubsection{Patient-Led Research
Collaborative}\label{patient-led-research-collaborative}

The Patient-Led Research Collaborative was formed in May 2020 by people
with Long COVID who met in an online support group
\citep{mccorkell2021}. Patients designed a survey, drawing on their own
and their communities' experience to compile the list of symptoms,
translated it into eight languages, and analysed responses from 3,762
people \citep{davis2020preprint, davis2021}. Fatigue, post-exertional
malaise (worsening after effort) and cognitive dysfunction were the most
common symptoms, reported by 98.3\%, 89.0\% and 85.1\% of respondents,
respectively. The study had ethics approval and the authors reported its
main limitation, a sample drawn from active support-group members. The
collaborative published the first report on Long COVID
\citep{fitzgerald2024}, and its members, all patients, were the authors.
In 2022, a panel of 15 patient-researchers allocated USD 4.8 million in
research grants to ten projects \citep{plrc_fund}. We found no
documentation of whether survey respondents received results directly,
or of public access to the dataset.

\emph{Lesson:} when people who live with a condition lead the research,
questions that conventional research has not yet asked come to light
quickly. Credibility came from rigorous methods and peer-reviewed
publication, not from patient status alone.

\subsubsection{Youth-GEMS}\label{youth-gems}

Youth-GEMS is a five-year project funded by Horizon Europe with 18
partners in ten countries, studying how genetic and environmental
factors shape mental health in people aged 12 to 24 \citep{prbb2023}.
Its research questions and design were set by the consortium. Young
Experts, aged 18 to 25 with lived experience of mental health problems,
are recruited with a partner youth organisation
\citep{youthgems_expert2025}. They shaped the wording, tone and visuals
of the project's mobile app and identified problems that technical teams
had not considered, but they were explicitly not involved in building or
validating the artificial-intelligence model, which remained with
technical experts and psychiatrists \citep{youthgems_decisions2025}. A
framework for regular involvement, data protection and safe spaces for
sharing experience was established, and young people's feedback is used
to assess the quality of engagement \citep{youthgems_earma2024}. Young
Experts receive training and vouchers rather than pay
\citep{youthgems_expert2025}. A clinical study of about 1,000 young
people in Europe and 1,000 in Australia was under way at the time of
writing, and no results had been published \citep{youthgems_research}.

\emph{Lesson:} institution-led projects tend to limit citizen roles to
consultation on selected components. Stating openly where involvement
stops, and why, is more credible than implying full partnership.

\subsubsection{LaporCovid-19}\label{laporcovid-19}

LaporCovid-19 was founded in March 2020 by a coalition of public health
experts, scientists, journalists and activists in response to the lack
of reliable COVID-19 data in Indonesia
\citep{laporcovid2021, rinaldi2024}. Citizens reported cases, deaths,
and experiences through a chatbot on WhatsApp and Telegram. From
September 2020, about 35 volunteers maintained a database of
health-worker deaths, searching obituaries in online news and social
media and collecting testimonies from families and colleagues
\citep{ekawati2022}. Using these data, an analysis with ethics approval
from the Ministry of Health's research ethics committee documented 1,545
health-worker deaths and estimated that health workers died at 4.92
times the rate of the general population, at a time when no official
count of health-worker infections or deaths existed \citep{ekawati2022}.
Aggregated data were made publicly available. The group advocated for
compensation for families of health workers, and civil-society pressure
contributed to the withdrawal of a self-funded vaccination plan
\citep{laporcovid2021}. Its relationship with government was contested,
as its data directly challenged official figures \citep{laporcovid2021}.
The organisation behind it continued after the emergency phase, renamed
LaporSehat and LaporIklim \citep{wargaberdaya2024, wargaberdaya_web},
but we found no documentation of whether the health-worker death
database is still maintained or of measures to protect volunteers.

\emph{Lesson:} citizen data can document what official statistics leave
out. A critical stance toward government protects independence but may
limit uptake.

\subsubsection{Community-based surveillance by the Indonesian Red
Cross}\label{community-based-surveillance-by-the-indonesian-red-cross}

The Indonesian Red Cross (\emph{Palang Merah Indonesia}) ran
community-based surveillance in eight villages across four districts,
including Boyolali in Central Java, as part of an international epidemic
preparedness programme \citep{ifrc2021}. The system was designed by the
Red Cross with community and government stakeholders. Volunteers
recruited from local structures, such as community leaders, midwives and
teachers, were trained to recognise and report signs of priority
diseases. Between September 2019 and September 2020, they raised 55
alerts, of which 48 were verified as true, and every alert received a
response within 24 hours. Alerts were fed into the national
early-warning system, and local governments took ownership; Boyolali
prioritised the programme in its local budget \citep{ifrc2021}. Data
were protected by recording only initials and sending numeric codes by
text message. Volunteers reported to a protocol designed by others and
had no role in analysis. Volunteers were unpaid, which the report
identifies as a challenge for retention. Of 234 volunteers trained in
surveillance, 171 (73\%) remained active.

\emph{Lesson:} building citizen reporting into existing government
systems makes it last, at the cost of keeping citizens at the reporting
stage. Volunteer retention needs planning and resources.

\subsubsection{Crisis Text Line data
sharing}\label{crisis-text-line-data-sharing}

Crisis Text Line is a free text-message crisis service in the United
States, staffed by trained volunteer counsellors. It ran a programme
sharing anonymised conversation data with academic researchers under
published guidelines for non-commercial use. It also created a
for-profit subsidiary, Loris.ai, which was licensed Crisis Text Line
data to train customer-service software \citep{eysenbach2025}. People
texting in crisis had agreed only to a terms-of-service link. After the
arrangement was reported in the press in January 2022 and a commissioner
of the US Federal Communications Commission called for it to stop,
Crisis Text Line ended the relationship within three days and asked
Loris.ai to delete the data. A former volunteer had raised concerns
about consent and conflict of interest, and the journal that had
published the original data-sharing guidelines issued a correction
clarifying the authors' ties to the commercial and non-commercial
entities \citep{eysenbach2025}. Texters had no role in the research
beyond being its data subjects.

\emph{Lesson:} consent given in a moment of crisis, through a
terms-of-service link, does not cover a later commercial use. A change
of purpose requires new consent and new ethical review, and ties to
commercial beneficiaries must be disclosed.

\subsubsection{Samaritans Radar}\label{samaritans-radar}

Samaritans Radar was an app launched by the UK suicide-prevention
charity Samaritans on 29 October 2014 \citep{responsibledata_radar}.
Based on university research, it scanned the tweets of accounts that an
app user followed for language associated with suicidal intent and
emailed the user an alert with suggested actions. The people whose
tweets were scanned were not told. Monitoring was on by default; people
could ask to be excluded only if they knew the app existed. Many in the
UK mental health community objected that tweets which were technically
public were not expected to be monitored, that people might stop
sharing, and that stalkers or bullies could use the alerts to find
people at their most vulnerable \citep{responsibledata_radar}.
Samaritans suspended the app nine days after launch
\citep{irishtimes2014} and closed it permanently in March 2015, deleting
all associated data. The charity stated that it should have consulted
more widely, and it later ran a public consultation on digital tools
\citep{responsibledata_radar}.

\emph{Lesson:} a post being technically public does not mean its author
has agreed to be monitored. Tools that use mental-health data should be
designed with the people being watched, not only the people watching,
and should require opt-in.

\subsection{Cross-case comparison}\label{cross-case-comparison}

Table~\ref{tbl-case-comparison} shows the scores for all eight cases. No
item was met in six or more cases, so no practice qualified as done well
across cases. Twelve of the eighteen items were met in two or fewer
cases.

\begin{landscape}

\begin{longtable}[]{@{}
  >{\raggedright\arraybackslash}p{(\linewidth - 18\tabcolsep) * \real{0.3100}}
  >{\centering\arraybackslash}p{(\linewidth - 18\tabcolsep) * \real{0.0800}}
  >{\centering\arraybackslash}p{(\linewidth - 18\tabcolsep) * \real{0.0800}}
  >{\centering\arraybackslash}p{(\linewidth - 18\tabcolsep) * \real{0.0800}}
  >{\centering\arraybackslash}p{(\linewidth - 18\tabcolsep) * \real{0.0800}}
  >{\centering\arraybackslash}p{(\linewidth - 18\tabcolsep) * \real{0.0800}}
  >{\centering\arraybackslash}p{(\linewidth - 18\tabcolsep) * \real{0.0800}}
  >{\centering\arraybackslash}p{(\linewidth - 18\tabcolsep) * \real{0.0800}}
  >{\centering\arraybackslash}p{(\linewidth - 18\tabcolsep) * \real{0.0800}}
  >{\centering\arraybackslash}p{(\linewidth - 18\tabcolsep) * \real{0.0500}}@{}}

\caption{\label{tbl-case-comparison}Scores by case}

\tabularnewline

\toprule\noalign{}
\begin{minipage}[b]{\linewidth}\raggedright
Item
\end{minipage} & \begin{minipage}[b]{\linewidth}\centering
CoAct
\end{minipage} & \begin{minipage}[b]{\linewidth}\centering
HGN
\end{minipage} & \begin{minipage}[b]{\linewidth}\centering
PLRC
\end{minipage} & \begin{minipage}[b]{\linewidth}\centering
YG
\end{minipage} & \begin{minipage}[b]{\linewidth}\centering
LC
\end{minipage} & \begin{minipage}[b]{\linewidth}\centering
CBS
\end{minipage} & \begin{minipage}[b]{\linewidth}\centering
CTL
\end{minipage} & \begin{minipage}[b]{\linewidth}\centering
SR
\end{minipage} & \begin{minipage}[b]{\linewidth}\centering
\texttt{Y} n/8
\end{minipage} \\
\midrule\noalign{}
\endhead
\bottomrule\noalign{}
\endlastfoot
\textbf{A1} Choosing the question & \texttt{P} & \texttt{N} & \texttt{Y}
& \texttt{N} & \texttt{Y} & \texttt{N} & \texttt{N} & \texttt{N} & 2 \\
\textbf{A2} Designing methods and tools & \texttt{Y} & \texttt{NR} &
\texttt{Y} & \texttt{P} & \texttt{Y} & \texttt{P} & \texttt{N} &
\texttt{N} & 3 \\
\textbf{A3} Collecting data & \texttt{P} & \texttt{P} & \texttt{Y} &
\texttt{NR} & \texttt{Y} & \texttt{P} & \texttt{N} & \texttt{N} & 2 \\
\textbf{A4} Analysing and interpreting data & \texttt{P} & \texttt{NR} &
\texttt{Y} & \texttt{N} & \texttt{Y} & \texttt{N} & \texttt{N} &
\texttt{N} & 2 \\
\textbf{A5} Writing and sharing results & \texttt{P} & \texttt{NR} &
\texttt{Y} & \texttt{P} & \texttt{Y} & \texttt{N} & \texttt{N} &
\texttt{N} & 2 \\
\textbf{A6} Acting on findings & \texttt{Y} & \texttt{NR} & \texttt{Y} &
\texttt{NR} & \texttt{Y} & \texttt{P} & \texttt{N} & \texttt{P} & 3 \\
\textbf{B1} Both citizens and science benefited & \texttt{Y} &
\texttt{Y} & \texttt{Y} & \texttt{P} & \texttt{Y} & \texttt{P} &
\texttt{P} & \texttt{P} & 4 \\
\textbf{B2} Participants got results back & \texttt{P} & \texttt{Y} &
\texttt{P} & \texttt{NR} & \texttt{NR} & \texttt{P} & \texttt{NR} &
\texttt{N} & 1 \\
\textbf{B3} Data and results are open & \texttt{P} & \texttt{NR} &
\texttt{P} & \texttt{NR} & \texttt{Y} & \texttt{P} & \texttt{P} &
\texttt{NR} & 1 \\
\textbf{B4} Contributions acknowledged in a form contributors chose &
\texttt{P} & \texttt{NR} & \texttt{P} & \texttt{P} & \texttt{P} &
\texttt{P} & \texttt{NR} & \texttt{NR} & 0 \\
\textbf{B5} Data quality checked, limitations stated & \texttt{P} &
\texttt{NR} & \texttt{Y} & \texttt{NR} & \texttt{Y} & \texttt{Y} &
\texttt{NR} & \texttt{NR} & 3 \\
\textbf{B6} Ethics and data protection in place & \texttt{P} &
\texttt{P} & \texttt{Y} & \texttt{P} & \texttt{Y} & \texttt{P} &
\texttt{P} & \texttt{N} & 2 \\
\textbf{B7} Project evaluated itself & \texttt{P} & \texttt{NR} &
\texttt{NR} & \texttt{P} & \texttt{NR} & \texttt{Y} & \texttt{NR} &
\texttt{P} & 1 \\
\textbf{C1} Citizens changed the research question & \texttt{P} &
\texttt{N} & \texttt{Y} & \texttt{N} & \texttt{Y} & \texttt{N} &
\texttt{N} & \texttt{N} & 2 \\
\textbf{C2} Findings came before the official or usual route &
\texttt{NR} & \texttt{NR} & \texttt{Y} & \texttt{NR} & \texttt{Y} &
\texttt{P} & \texttt{NR} & \texttt{NR} & 2 \\
\textbf{C3} Findings were trusted and used & \texttt{P} & \texttt{NR} &
\texttt{P} & \texttt{NR} & \texttt{P} & \texttt{Y} & \texttt{N} &
\texttt{N} & 1 \\
\textbf{D1} Activity lasted, or ended as planned & \texttt{P} &
\texttt{Y} & \texttt{Y} & \texttt{Y} & \texttt{P} & \texttt{Y} &
\texttt{N} & \texttt{N} & 4 \\
\textbf{D2} Volunteers and participants protected from harm & \texttt{Y}
& \texttt{NR} & \texttt{NR} & \texttt{Y} & \texttt{NR} & \texttt{Y} &
\texttt{P} & \texttt{P} & 3 \\

\end{longtable}

\textbf{\emph{Note.}}\\
\texttt{Y} met; \texttt{P} partly met; \texttt{N} not met; \texttt{NR}
not reported.\\
HGN, HoeGekIsNL; PLRC, Patient-Led Research Collaborative; YG,
Youth-GEMS; LC, LaporCovid-19; CBS, Community-based surveillance
(Indonesian Red Cross); CTL, Crisis Text Line; SR, Samaritans Radar.
Provisional, pending second coder.

\end{landscape}

\subsubsection{What changed because citizens took
part}\label{what-changed-because-citizens-took-part}

Citizens changed the research question in two cases. In both, the
initiative was founded by the group itself: patients defined Long COVID
as an object of study \citep{mccorkell2021, fitzgerald2024}, and the
LaporCovid-19 coalition and its volunteers counted health-worker deaths
that officials were not counting \citep{ekawati2022}. The same two cases
produced findings before official or conventional sources. Documented
use of findings by decision-makers was rarer, with only community-based
surveillance met this item, through integration into the national
early-warning system and local budgets \citep{ifrc2021}. In three other
cases, findings attracted attention or were presented to authorities
without documented adoption. In both failure cases, the activity was
discredited and ended.

\subsubsection{Where citizens took part}\label{where-citizens-took-part}

Design was the stage at which citizens most often had a leading role
(three of eight cases), followed by acting on findings (three of eight).
The founding group led every stage only in the two founder-led cases. In
LaporCovid-19, lay members of the public mainly reported data while the
coalition held the leading roles \citep{rinaldi2024, ekawati2022}. In
institution-led cases, citizen roles were concentrated in data
collection and consultation on specific components, such as the wording
of an app \citep{youthgems_decisions2025} or the reporting of alerts
\citep{ifrc2021}. In the two failure cases, people whose data were used
had no research role at all.

\subsubsection{How projects were run}\label{how-projects-were-run}

Three items were consistently lacking. Results were returned to
participants in a usable form in one case, HoeGekIsNL, through
personalised reports \citep{vanderkrieke2016, mindwise2016}. No case
documented that contributors could choose how they were credited; five
cases acknowledged contributors in some form, and three did not report
how. A formal evaluation was published in one case \citep{ifrc2021}.
Openness of both data and results was documented in one case
\citep{ekawati2022}, and ethics review together with a consent approach
in two \citep{ekawati2022, davis2020preprint}. Across these lacking
items, most cases scored \texttt{NR} rather than \texttt{N}: the sources
did not say whether the practice took place. The exception was consent
in the failure cases, where sources documented its absence or inadequacy
\citep{eysenbach2025, responsibledata_radar}.

\subsubsection{Durability and harm}\label{durability-and-harm}

Four cases were still running or had ended as planned, and two were
partly durable: CoAct, and LaporCovid-19, whose organisation continued
under new names while its original data activity is undocumented. The
two failures ended under public pressure. Protection from harm was
documented in three cases, through anonymous participation
\citep{casanovas2026}, safe spaces for young people
\citep{youthgems_earma2024}, and risk-aware protocols for volunteers
\citep{ifrc2021}. In both failures, the harm was a breach of trust in
how mental-health data were used, partly addressed by stopping the
activity and deleting data.

\subsubsection{Comparing successes and
failures}\label{comparing-successes-and-failures}

The successful and partly successful cases shared two features that the
failures lacked: the people whose experience was being studied had some
role in the research, and their data were used for the purpose they had
agreed to. Both failures began with an organisation acting on behalf of
people in distress without asking them, and both ended when those
people, or volunteers speaking for them, objected publicly.

\section{Discussion}\label{sec-discussion}

Across the eight cases, citizens held leading roles at every stage only
where the group that founded the initiative kept contributing to the
work, as in the Patient-Led Research Collaborative and LaporCovid-19. In
the institution-led cases, citizen involvement was mostly instrumental,
limited to tasks that the lead institution had already defined, such as
data collection or consultation on selected components. In the two
failures, the people whose data were used had no research role at all.
This is consistent with the systematic review of mental-health citizen
science, in which citizens contributed mostly through data collection
\citep{todowede2023}. Citizen leadership was rare in these cases. Where
it occurred, it shaped the research question and produced evidence
before official or conventional sources did.

Citizen involvement changed the research more readily than it changed
decisions. In the sources we reviewed, documented use of findings by
decision-makers was scarce, even in the two founder-led cases.
LaporCovid-19's advocacy contributed to the withdrawal of a self-funded
vaccination plan, but its data challenged official figures
\citep{laporcovid2021}. Findings from the Patient-Led Research
Collaborative drew attention, but we found no documented route into
decisions. The only case with documented use was the Indonesian Red
Cross community-based surveillance, whose alerts were fed into the
national early-warning system \citep{ifrc2021}. Its volunteers, however,
only reported to a protocol designed by others. These cases suggest a
trade-off, where working inside government systems kept citizens at the
reporting stage, while independent leadership by the founding group left
its relationship with government contested. No case clearly combined
citizen leadership with documented use of findings. The three conditions
below are synthesized from the cases and have not been tested. They bear
on a practical problem: how can citizen leadership over the research
question be combined with the institutional routes through which
findings become policy or programmes?

The first condition is a pathway from findings to decisions, agreed at
the design stage with a named institutional collaborator. In the Red
Cross case, the citizen-contributed alerts were fed into the national
early-warning system, and local governments took ownership to resolve
the issue \citep{ifrc2021}. On the other hand, CoAct produced 14
recommendations and public commitments from a city councillor and the
director of the National Mental Health Pact
\citep{ub2022assembly, bcn2023plan}. We could trace no adoption, which
suggests that commitments need to be recorded and followed up. In
LaporCovid-19, the health-worker database was built with professional
associations, a news organisation and a research unit, and the analysis
had Ministry of Health ethics approval \citep{ekawati2022}. This is a
single observation, but it shows that a civil-society database can gain
institutional standing without being run by an institution. Wider
evidence points the same way, with caution. A systematic review of 145
studies found that collaboration and relationships between researchers
and policymakers were among the most often reported facilitators of
evidence use \citep{oliver2014}. Poor access to relevant research, and
research arriving too late, were the main barriers. A scoping review of
integrated knowledge translation, in which research is done in
partnership with decision-makers from the outset, found that most
evaluated partnerships reported at least one positive outcome
\citep{gagliardi2016}. It is not known which partnership strategies lead
to which outcomes. A benchmark outside our eight cases is the NIH
RECOVER programme on Long COVID. Its adult-cohort protocol states that
patient representatives give input on all study processes, and some were
recruited from advocacy groups including the Patient-Led Research
Collaborative \citep{horwitz2023}. Representatives have described the
partnership from their side \citep{kim2023}. A survey of RECOVER's
clinical-trial protocol working groups found that most patient,
caregiver and community representatives felt their input was valued, but
gaps in role clarity, communication and integration of feedback
persisted \citep{leverty2026}. Embedding patients in a research
programme is not the same as institutions using findings that patients
generated (item C3). A news feature in \emph{Nature} reported that the
survey had been mentioned in some 60 policy statements
\citep{fairbank2024}. Overall, the benchmark suggests that influence on
research design and on policy discussion is achievable, while
decision-makers' use of findings is harder to show.

The second condition is legitimacy, and both failures broke down at
consent. In the Crisis Text Line case, data given in crisis under a
terms-of-service link were later licensed for a commercial purpose
\citep{eysenbach2025}. In Samaritans Radar, people were monitored
without knowing it, under an opt-out design
\citep{responsibledata_radar}. Both activities were discredited and
ended under public pressure. A similar loss of public trust ended
England's national plan to share general-practice records. That plan
proceeded without what has been called a ``social licence''
\citep{carter2015}. Positive approaches exist. The MindKind study asked
young people in India, South Africa and the United Kingdom how their
mental-health data should be governed \citep{mindkind2022}. It found
that they held strong preferences about who should access their data.
They were also willing to help design more open data-sharing
arrangements. In Indonesia, personal data protection is regulated by Law
27/2022 \citep{uupdp2022}, and the Health Law of 2023 contains
provisions on health information systems \citep{uu17_2023}. Although we
do not analyse their provisions here, these cases suggest that consent
for mental-health data should be specific to each study and recorded so
that it can be checked. It should be renewed when the purpose changes,
and opt-in by default. These requirements are developed as principles in
Paper 2 of this series.

The third condition is credibility, which depends on being able to see
how a project was run. Most of the consistently lacking items were
unreported rather than shown to be absent. From public sources,
therefore, we cannot tell whether a project omitted a practice or simply
did not write about it. The same gap was found in the mental-health
citizen science literature \citep{todowede2023}. Reporting guidance
already exists. The GRIPP2 checklists describe how to report patient and
public involvement \citep{staniszewska2017}. The systematic review
proposed mental-health-specific reporting recommendations for each ECSA
principle \citep{todowede2023}. Our 18 items could serve as a
complementary checklist, allowing projects to state explicitly which
practices they followed, including whether they evaluated themselves.
Reporting also needs to say whose experience is represented. Typologies
of participation, from Arnstein's ladder \citep{arnstein1969} to
Haklay's levels \citep{haklay2013}, rank projects by how much power
citizens hold. They do not distinguish which citizens hold it. In
LaporCovid-19, the leading roles were held by a coalition of experts,
journalists and activists, while lay members of the public mainly
reported data \citep{rinaldi2024, ekawati2022}. The Patient-Led Research
Collaborative was led by people with Long COVID who describe themselves
as patient-researchers \citep{mccorkell2021, fitzgerald2024}.
Participation dominated by a small or well-connected group risks the
``elite capture'' described in lived-experience research
\citep{okoroji2023}. The people most affected by distress, who may have
the least time, energy and confidence to take part, remain unheard.
Future frameworks should record not only the level of participation but
also whose experience is represented at each level.

\subsection{Recommendations for future research and
practice}\label{recommendations-for-future-research-and-practice}

Drawing on the reviewed cases, we propose a design sequence for
citizen-centred research on stress, depression, anxiety, burnout and
resilience. The sequence begins with citizens helping to formulate and
sharpen the research question, so that the study addresses what matters
to the people affected. In the two founder-led cases, this produced
evidence that official or conventional sources had not, such as the
health-worker deaths that no official source was counting
\citep{ekawati2022}. In the two failures, by contrast, organisations
acted on behalf of people in distress without asking them. In a
nine-country study, Indonesian respondents identified community support,
including mental support, as a way they wanted to contribute
\citep{tan2024}. Before data collection, projects should obtain opt-in,
study-specific consent. Where feasible, they should also agree with a
named institutional counterpart how the findings might be used. Results
should then be returned to participants in a usable form. The study
should be evaluated and reported against the 18 items, so that readers
can judge how it was run and whose experience it represents. Indonesia
has a channel to build on in lay mental health cadres. Evaluations of
their training, however, have not measured outcomes for patients or
communities, and support for volunteers appears uneven
\citep{susanti2024, sitorukmi2025}. Within this series, Paper 2 develops
the consent and reporting lessons into principles for research-ready
health data. Paper 3 assesses an open-source research platform against
those principles. Paper 4 is a baseline study of depression, anxiety,
burnout, resilience and functioning among adults in three Indonesian
cities. Beyond the series, the scoring framework needs validation, and
interviews with project teams could show whether unreported practices
took place. A pilot mental-health study in Indonesia, with a pre-agreed
pathway to a named institution, would test whether the sequence makes
findings actionable.

\subsection{Strengths and limitations}\label{strengths-and-limitations}

The main limitation of this study is its reliance on public documents.
Practices that were not reported could not be scored, so the many
\texttt{NR} scores more likely reflect gaps in reporting than gaps in
practice. The cases were selected purposively and are not representative
of citizen science in mental health, so the patterns described here
should be read as hypotheses and not as frequencies. Several cases also
rested on organisational reports, project websites or news coverage,
which may present projects favourably or selectively. This applies in
particular to the D1 score for LaporCovid-19, which rests on
organisational sources. That score was revised from not reported to
partly met after a follow-up search on 16 November 2026. The scoring
framework is new and has not been validated beyond this review, and
agreement between coders is reported only as a percentage of cells
agreed. Defining the unit of analysis and who counts as a citizen also
required judgement, particularly for the failure cases, in which the
people whose data were used had no research role. In addition, the lead
author founded a company that develops a research platform evaluated
later in this series; the steps taken to limit bias are described in the
Methods and Declarations. Finally, we found no Indonesian citizen
science project focused on mental health. This may reflect a true
absence or the limits of our search, and the Indonesian recommendations
therefore draw on benchmarks from other fields and countries.

Notwithstanding these limitations, the review has several merits. All
cases were scored by two coders on a common framework, with rules fixed
before scoring. Failure cases were included and scored on the same terms
as successes, which limits the tendency to describe only what worked.
The review also adds evidence from a middle-income country, addressing a
gap in the existing literature \citep{todowede2023}.

\section{Conclusion}\label{sec-conclusion}

Citizen involvement in health research can change which questions are
asked and produce evidence before conventional routes do, particularly
when the people affected lead the work. The weak points are in how
projects are run and reported, where feedback to participants, credit,
evaluation and consent are rarely documented. Citizen-centred research
on stress, depression, anxiety, burnout and resilience in Indonesia
should be designed around these weak points from the start, and should
report them explicitly so that others can learn from both its successes
and its failures.

\section*{References}\label{references}
\addcontentsline{toc}{section}{References}

\renewcommand{\bibsection}{}
\bibliography{../references.bib}





\end{document}
